\documentclass[10pt,letterpaper]{article}
\usepackage[T1]{fontenc}
\usepackage[utf8]{inputenc}
\usepackage[spanish,es-tabla]{babel}
\usepackage{csquotes}
\usepackage{mathpazo}
\usepackage[scaled=0.9]{helvet}
\usepackage{courier}
\usepackage{calc}
\usepackage[top=2cm, bottom=2cm, left=2cm, right=2cm]{geometry}
\usepackage{xcolor}
\usepackage{booktabs}
\usepackage{longtable}
\usepackage{array}
\usepackage{ragged2e}
\usepackage{xurl}
\usepackage[strings]{underscore}
\usepackage{fancyhdr}
\usepackage{sectsty}
\usepackage{tocloft}
\usepackage[colorlinks=true,linkcolor=blue,urlcolor=blue,citecolor=blue,breaklinks=true]{hyperref}
\usepackage{enumitem}

\definecolor{uncpblue}{HTML}{1A237E}
\allsectionsfont{\color{uncpblue}}

\renewcommand{\cftsecleader}{\cftdotfill{\cftsecdotsep}}
\cftpagenumbersoff{section}
\cftpagenumbersoff{subsection}
\cftpagenumbersoff{subsubsection}
\renewcommand{\cftsecfont}{\normalfont}
\renewcommand{\cftsubsecfont}{\normalfont}
\renewcommand{\cftsubsubsecfont}{\normalfont}
\setlength{\parindent}{0pt}
\setlength{\parskip}{4pt}
\setlength{\tabcolsep}{3pt}
\renewcommand{\arraystretch}{0.85}

\pagestyle{fancy}
\fancyhf{}
\fancyhead[L]{\small\textcolor{uncpblue}{\textbf{SGSI-UNCP}}}
\fancyhead[R]{\small\textcolor{uncpblue}{ACT-SGSI-01}}
\fancyfoot[C]{\thepage}
\renewcommand{\headrulewidth}{0.4pt}
\renewcommand{\footrulewidth}{0.4pt}
\setlength{\headheight}{14pt}

\setlist{nosep,leftmargin=*,after=\vspace{2pt}}
\setlength{\emergencystretch}{2em}

\newcommand{\makecover}{
\begin{titlepage}
\centering
\vspace*{2cm}
{\Huge\bfseries\color{uncpblue} SGSI-UNCP\par}
\vspace{0.7cm}
{\LARGE Acta de Compromiso de la Alta Dirección con el SGSI\par}
\vspace{0.5cm}
{\large Documento Base del Sistema de Gestión\\de Seguridad de la Información\par}
\vspace{1.5cm}
{\small Universidad Nacional del Centro del Perú\par}
\vfill
\end{titlepage}
}

\begin{document}

\makecover
\setcounter{tocdepth}{2}
\RaggedRight
\tableofcontents
\newpage

\RaggedRight
\section{Acta de Compromiso de la Alta Dirección con el
SGSI}\label{acta-de-compromiso-de-la-alta-direcciuxf3n-con-el-sgsi}

\textbf{Organización:} Universidad Nacional del Centro del Perú (UNCP)
\textbf{Aprobado por:} Rector / Presidente del Comité de Gobierno
Digital \textbf{Norma:} ISO/IEC 27001:2022 (Cláusula 5.1)

\begin{center}\rule{0.5\linewidth}{0.5pt}\end{center}

\subsection{Declaración de
Intención}\label{declaraciuxf3n-de-intenciuxf3n}

La Alta Dirección de la Universidad Nacional del Centro del Perú (UNCP),
en su calidad de máxima autoridad y en el marco de sus competencias
establecidas en la Ley Universitaria N.° 30220 y el Estatuto
Institucional, reconoce que la información académica, administrativa e
investigativa constituye un activo estratégico fundamental para el
cumplimiento de la misión universitaria.

En concordancia con el Plan de Gobierno y Transformación Digital (PGTD
2026-2030) y las disposiciones del Decreto Legislativo N.° 1412 (Ley de
Gobierno Digital), manifestamos nuestro compromiso incondicional con el
establecimiento, implementación, mantenimiento y mejora continua del
Sistema de Gestión de Seguridad de la Información (SGSI) basado en la
norma internacional ISO/IEC 27001:2022.

\subsection{Compromisos Específicos}\label{compromisos-especuxedficos}

Para asegurar la eficacia del SGSI y en cumplimiento de la Cláusula 5.1
de la ISO/IEC 27001:2022, la Alta Dirección asume los siguientes
compromisos:

\begin{enumerate}
\def\labelenumi{\arabic{enumi}.}
\item
  \textbf{Liderazgo y Ejemplo:} Actuar como promotores visibles de la
  cultura de seguridad de la información en toda la comunidad
  universitaria, participando activamente en las actividades de
  sensibilización y difusión del SGSI.
\item
  \textbf{Asignación de Recursos:} Dotar al SGSI del presupuesto,
  personal calificado, infraestructura tecnológica y tiempo necesario
  para su implantación y operación continua, conforme al proyecto
  PGTD-01 (S/ 850,000).
\item
  \textbf{Alineamiento Estratégico:} Asegurar que los objetivos de
  seguridad de la información estén integrados en los procesos de
  planeamiento institucional y alineados con los Objetivos Estratégicos
  Institucionales (OEI) del PEI 2026-2030 y los Objetivos de Gobierno y
  Transformación Digital (OGTD) del PGTD.
\item
  \textbf{Revisión por la Dirección:} Realizar revisiones periódicas del
  SGSI al menos una vez al año o cuando ocurran cambios significativos,
  evaluando su conveniencia, adecuación, eficacia y eficiencia, de
  acuerdo con la Cláusula 9.3 de la ISO/IEC 27001:2022.
\item
  \textbf{Comunicación Interna y Externa:} Promover la importancia de
  una gestión de seguridad eficaz mediante canales formales de
  comunicación, asegurando que todas las partes interesadas comprendan
  los requisitos del SGSI y su rol en el cumplimiento de los mismos.
\item
  \textbf{Mejora Continua:} Impulsar la adopción del ciclo PHVA
  (Planificar-Hacer-Verificar-Actuar) como metodología para la mejora
  continua del SGSI, fomentando la identificación y corrección proactiva
  de desviaciones.
\item
  \textbf{Cumplimiento Normativo:} Garantizar el cumplimiento de la Ley
  N.° 29733 (Ley de Protección de Datos Personales), el D.S. N.°
  029-2021-PCM (Reglamento de la Ley de Gobierno Digital), el Marco de
  Confianza Digital y demás normativa aplicable en materia de seguridad
  digital y protección de datos.
\end{enumerate}

\subsection{Designación de Autoridad}\label{designaciuxf3n-de-autoridad}

En cumplimiento del Artículo 5 del D.S. N.° 029-2021-PCM, se designa y
ratifica al:

\begin{itemize}
\tightlist
\item
  \textbf{Director(a) General de Administración} como Líder de Gobierno
  y Transformación Digital, responsable de la supervisión estratégica
  del SGSI y de la ejecución de los proyectos del PGTD.
\item
  \textbf{Oficial de Seguridad y Confianza Digital} como responsable
  operativo del SGSI, con plena autoridad técnica para la implementación
  de controles, gestión de incidentes y administración de riesgos de
  seguridad de la información.
\end{itemize}

Ambos cargos reportarán directamente al Comité de Gobierno Digital y al
Rectorado sobre el estado y desempeño del SGSI.

\subsection{Incumplimiento}\label{incumplimiento}

El incumplimiento de los compromisos establecidos en la presente acta o
de las políticas y controles derivados del SGSI será puesto en
conocimiento del Comité de Gobierno Digital para la adopción de medidas
correctivas, sin perjuicio de las responsabilidades administrativas,
civiles o penales que pudieran corresponder conforme a la normativa
universitaria y la legislación nacional aplicable.

\begin{center}\rule{0.5\linewidth}{0.5pt}\end{center}

\textbf{Firmas de la Alta Dirección:}

\begin{center}\rule{0.5\linewidth}{0.5pt}\end{center}

\textbf{Rector de la UNCP} Presidente del Comité de Gobierno Digital

\begin{center}\rule{0.5\linewidth}{0.5pt}\end{center}

\textbf{Vicerrector Académico}

\begin{center}\rule{0.5\linewidth}{0.5pt}\end{center}

\textbf{Vicerrector de Investigación}

\begin{center}\rule{0.5\linewidth}{0.5pt}\end{center}

\textbf{Director(a) General de Administración} Líder de Gobierno y
Transformación Digital

\begin{center}\rule{0.5\linewidth}{0.5pt}\end{center}

\textbf{Oficial de Seguridad y Confianza Digital}

\begin{center}\rule{0.5\linewidth}{0.5pt}\end{center}

\emph{Documento aprobado en Sesión del Comité de Gobierno Digital, en
conformidad con el D.S. N.° 029-2021-PCM y la R.S. N.°
005-2018-PCM-SEGDI.}

\end{document}
